\section{Introduction}\label{sec:introduction}
A linear isoperimetric inequality is a large-scale expression of negative curvature: for finitely presented groups, it characterizes word hyperbolicity. Local comparison geometry is more restrictive. The geometric realization question asks whether every hyperbolic group admits a geometric action on a proper $\CAT(0)$ space, or on a proper $\CAT(-1)$ space. We construct a group for which neither action exists.

Throughout the paper, a \emph{geometric action} is a proper, cocompact action by isometries, and metric models of finite complexes carry geodesic length metrics inducing the given complex topology.

\begin{theorem}\label{thm:main}
There exist a finite connected two-dimensional aspherical simplicial complex $K$ and a constant $C<\infty$ with the following properties.
\begin{enumerate}[label=\textup{(\roman*)}]
\item Every combinatorial edge loop of length $\ell$ in the universal cover of $K$ bounds a disk diagram with at most $C\ell$ triangular faces.
\item The group $\pi_1K$ admits no geometric action on a proper complete $\CAT(0)$ space.
\end{enumerate}
In particular, $\pi_1K$ is word-hyperbolic, and no finite simplicial complex homotopy equivalent to $K$ admits a locally $\CAT(0)$ geodesic length metric inducing its given topology. Consequently, no such model admits a compatible locally $\CAT(-1)$ geodesic length metric.
\end{theorem}

The proof is an existence argument and does not give an effective size for $K$. The filling constant may depend on the selected finite presentation.

\subsection{Context}
The realization problem goes back to the development of hyperbolic groups.
In his account of that development, Gromov recalls introducing linear
isoperimetry as a group-theoretic form of negative curvature and then
spending years seeking geometric realizations by negatively curved
spaces~\cite[p.~396]{RaussenSkau2010}.
The equivalence between linear filling and word hyperbolicity became
part of the foundational theory~\cite{Gromov1987}; see also
\cite[Theorem~III.$\Gamma$.2.6 and Remark~III.$\Gamma$.2.7(2)]{BH1999}.
The $\CAT(-1)$ realization question appears as Problem~H11 in
Baumslag--Myasnikov--Shpilrain's 1999 problem
list~\cite{BaumslagMyasnikovShpilrain1999}.
Both the $\CAT(-1)$ and $\CAT(0)$ questions appear as open problems
in Stark's 2025 survey~\cite[Open Problem~2.15]{Stark2025} and in the
March 2026 preprint of Duchesne--Simon~\cite[Question~W]{DuchesneSimon2026}.

There are several successful ways to obtain such geometric actions.
Finitely generated free groups act on trees and closed hyperbolic manifold groups act on
hyperbolic space. More generally, a finite complex carrying a locally
$\CAT(-1)$ or locally $\CAT(0)$ geodesic length metric inducing its topology has a universal cover
with the corresponding global curvature bound; its deck action is
geometric~\cite[Theorem~II.4.1]{BH1999}.
Cubulation provides another route to $\CAT(0)$ models.
For example, Ollivier--Wise prove that random groups in the density
model at density less than $1/6$ act freely and cocompactly on
$\CAT(0)$ cube complexes with probability tending to one as relator
length grows, with the number of generators fixed~\cite[Theorem~10.4]{OllivierWise2011}.
Thus random presentations can produce geometric $\CAT(0)$ actions as
well as hyperbolicity; the relations used here must enforce an
additional obstruction.

Dimension restrictions already yield obstructions.
Brady--Crisp~\cite[Theorem~2]{BradyCrisp2007} construct infinitely many
torsion-free hyperbolic groups of geometric dimension two which admit
no proper isometric action on a proper two-dimensional $\CAT(0)$
space. Their Theorem~1 gives a hyperbolic group of $\CAT(0)$ dimension
two and $\CAT(-1)$ dimension three. These conclusions rely on
constraints on translation lengths in two-dimensional targets.
Theorem~\ref{thm:main} excludes geometric $\CAT(0)$ actions in every
dimension. Its finite-model consequence also allows the entire
simplicial complex to change within its homotopy type: any compatible
locally $\CAT(0)$ geodesic length metric on such a model would recover a geometric
action by passing to the universal cover.

Random-presentation proofs of hyperbolicity combine counting for
bounded diagrams with a local-to-global argument.
Our treatment of repeated relator occurrences is related to the
multiplicity and constraint methods in
Ollivier~\cite[Section~2.2]{Ollivier2004} and
Antoniuk--Friedgut--\L uczak~\cite[Section~3]{AntoniukFriedgutLuczak2017}.
Within each finite presentation we first sample small labelled graphs,
then larger labelled graphs whose required small graphs have already
been selected. The same prerequisite can support several larger
graphs, so its probability cost must be charged only once.
Section~\ref{sec:diagrams} proves the resulting estimate directly.

The obstruction to actions uses the harmonic-energy framework of
Gromov~\cite[Sections~3.3 and~3.7]{Gromov2003} and
Izeki--Nayatani~\cite{IzekiNayatani2005}.
In this framework, expansion in a finite graph can force a fixed point
when the target satisfies a suitable nonlinear spectral inequality.
The target hypothesis matters: Izeki--Kondo--Nayatani obtain random
hyperbolic groups with fixed-point properties for target classes whose
tangent-cone invariant is bounded by a prescribed constant strictly
less than one~\cite[Corollary~3.8]{IzekiKondoNayatani2012}.
Toyoda~\cite[Theorem~1.4]{Toyoda2016} supplies such spectral bounds
for a prescribed finite family of geodesically complete $\CAT(0)$
spaces admitting proper cocompact actions, with constants depending
on that family.
Here a hypothetical target and its action can be chosen after the
finite presentation is known. We obtain the required uniformity from
bounds on every generator displacement and total-variation couplings
of table matches. Sturm's inductive
means~\cite[Definition~4.6 and Theorem~4.7, proof~(a)]{Sturm2003}
then turn those couplings into the averaging estimate used in the
energy argument.

\subsection{Structure of the proof}
The construction must satisfy two requirements: its relations must be sparse
enough to preserve hyperbolicity and asphericity, yet impose enough metric
constraints to exclude geometric actions on proper $\CAT(0)$ spaces.
We organize the relations using a fixed modification of the incidence graph
of the projective plane over $\mathbb F_2$.  Four point vertices are doubled.
At one of them we omit one incidence and keep its counterpart at the other
representative separate.  The remaining edges over each point form a small
connected graph, called a fan.  We call the resulting incidence graph
the core.  Edges receive signed generator labels,
and every closed walk in a selected labelled graph gives a relation.
Fans are selected first; a labelled copy of the whole graph is selected only
when its seven fan restrictions have already been selected.  We study
copies for which further four-letter relations close every nonbacktracking two-edge core
path to a square.  An independent collection of six-letter relations will
control individual generator displacements.  A labelled core copy together
with one chosen square completion for every nonbacktracking two-edge core path is a
\emph{match}.

We choose a sequence of such finite presentations, with progressively stronger
finite counting requirements.  At each stage, a finite list of statistical
tests is fixed before choosing the sampling parameters; those parameters then
remain fixed while the alphabet grows.  If infinitely many of their groups admitted
geometric actions on proper complete $\CAT(0)$ spaces, choose one
action for each of those groups.
The closed-walk relations realize the fixed graph by orbit points in every
chosen action.  After normalizing the average squared generator displacement,
products of these configurations have a limiting configuration in a complete
$\CAT(0)$ space.  The counting requirements force exact metric identities in
that limit, and a finite geometric contradiction excludes it.  Thus every
sufficiently late group in the chosen sequence admits no such action.  The argument uses the following three ingredients.

\begin{enumerate}
\item \textbf{A positive deficit in diagram counting.}
A fan shared by several copies of the whole graph supplies only one random
table entry, so its probability cost must be counted once.  At the same time,
a disk or sphere diagram may use the same relation many times.  We reconcile
these two counts by assigning multiplicity weights to distinct table entries.
Local incidence estimates at vertices then leave a positive multiple
of the diagram's area to be paid by its boundary.  This is first proved for
bounded fragments.  Planar separators extend it to all diagrams, yielding
both a linear filling inequality and asphericity in
Section~\ref{sec:diagrams}.

\item \textbf{Sampling that is uniform over later actions.}
The metric argument compares a bounded function of one or two edge labels
with a further edge label.  The relevant tuples need not have dense
projections in the sampled tables.  Instead, we couple several matches
so that the conditioning labels agree, each match is approximately
uniform among all matches, and the further labels are approximately independent
uniform generators.  Section~\ref{sec:sampling} obtains this coupling by
gluing copies along an enlarged conditioning tuple and proving concentration
of the numbers of extensions.  The enlargement is chosen by counting
exponents so that every further collection of label positions has a strictly positive
margin.  The conclusions concern total variation of finite probability laws.  They therefore apply to bounded tests from any action
chosen after the tables are revealed.  The common conditioning tuple may still depend on the further labels; the
averaging estimate allows this dependence.  The six-letter relations supply
the uniform displacement bound that makes these tests bounded.

\item \textbf{A rigid configuration and incompatible scalar products.}
Energy minimization and the coupling yield lower bounds for squared orbit
distances along short paths.  In the limit, the decoration squares force
squared two-step distances to equal two.  The six-cycles through three noncollinear point vertices then have a
common center, and their central triangles are flat
(Section~\ref{sec:energy}).  The doubled points give four unit vectors obtained
by doubling certain midpoints in the tangent cone at that center.  The line
incidences force two of these vectors to have nonnegative scalar product.
A three-field test involving the separated incidence edge forces the other
two vectors to be antipodal.  Splitting the parallel set of the resulting line exposes a flat
square in the transverse factor, which makes the first pair antipodal as well.
Their scalar product is consequently $-1$, the contradiction in
Section~\ref{sec:obstruction}.
\end{enumerate}

Section~\ref{sec:model} defines the finite graph and its random tables.
Section~\ref{sec:completion} makes the finite-test selection precise and
transfers one resulting finite classifying complex to the simplicial model
in Theorem~\ref{thm:main}.  Its filling constant may depend on that chosen
presentation.

\subsection{Conventions}
A disk diagram may have boundary spurs or edges collapsed to the target's $1$-skeleton; such features have zero area and can be removed. Its area is the number of its noncollapsed $2$-cells. Constants in filling estimates may depend on a fixed finite presentation. A statement holds \emph{with high probability} if its probability tends to one as the alphabet size tends to infinity, with all parameters in that statement fixed. Finite-test limits and the final diagonal selection will always be distinguished.

We use completed tangent cones and completed spaces of directions. The latter are $\CAT(1)$ and the former are $\CAT(0)$~\cite[Theorem~II.3.19]{BH1999}. For the tangent-cone construction in this generality, see also
Kleiner--Leeb~\cite[Section~2.1.3]{KleinerLeeb1997}, who credit
Nikolaev's tangent-cone theorem.
Subsequent products and limits need not be proper; properness is used only for the original geometric actions. Standard comparison, covering-space and topological results are cited with the hypotheses needed here. All estimates specific to the new construction are proved below.
